\documentclass[11pt,reqno]{amsart}

\usepackage[T1]{fontenc}
\usepackage[utf8]{inputenc}
\usepackage{lmodern}
\usepackage{amsmath,amssymb,amsthm}
\usepackage{booktabs}
\usepackage{array}
\usepackage[margin=1.1in]{geometry}
\usepackage[colorlinks=true,linkcolor=black,citecolor=black,urlcolor=blue]{hyperref}
\usepackage{microtype}

\theoremstyle{plain}
\newtheorem{theorem}{Theorem}
\newtheorem*{lemma}{Lemma}

\theoremstyle{remark}
\newtheorem*{verification}{Verification}
\newtheorem*{remark}{Remark}
\newtheorem*{rangeofstatement}{Range of the statement}

\numberwithin{equation}{section}

\newcommand{\fl}[1]{\left\lfloor #1 \right\rfloor}

\begin{document}

\title[An exact histogram for a quadratic staircase]
      {An Exact Histogram for a Quadratic Staircase}

\subjclass[2020]{Primary 11B57; Secondary 11A07, 11K31}
\keywords{Floor function, quadratic residues, exact multiplicities,
          Sturmian words, Beatty and step sequences,
          three-distance theorem, equidistribution}

\author{Mohamed Osman}
\address{Independent researcher}
\email{mohamedosmannn2999@gmail.com}
\urladdr{ORCID: 0009-0004-5912-999X}

\date{\today}

\begin{abstract}
For odd $n$ let
\[
  W_j \;=\; \fl{\frac{2(j+1)^2}{n}} - \fl{\frac{2j^2}{n}},
  \qquad j = 0,1,\dots,n-1 .
\]
We prove that $W_j \in \{0,1,2,3,4\}$ for every odd $n$ and every $j$ in this
range, and that the five values occur with multiplicities determined by the
single integer $A = \fl{(n+7)/8}$:
\[
  \#\{W{=}0\} = \#\{W{=}4\} = A, \qquad \#\{W{=}2\} = 2A-1, \qquad
  \#\{W{=}1\} = \#\{W{=}3\} = \tfrac{n+1}{2} - 2A .
\]
The proof is a tiling argument and carries no error term. This is worth
emphasising: a probabilistic model of the same count, assuming
$2j^2 \bmod n$ equidistributed, returns the same answer, and sums of this
shape normally carry an error of size $\sqrt{n}\log n$ or $n^{1/3}$. None
appears, and the residues $2j^2 \bmod n$ never enter the argument.

A companion compression is proved for the symmetric pairs (Theorems~3 and~4):
a pair collapses to a single element of $\{0,1,2\}$, and its census is again
governed by one integer.

We then prove the corresponding statement about the residues themselves, which
the histogram deliberately avoids: for odd $n \ge 51$ the sequence
$2j^2 \bmod n$ has exactly $2A$ interior local maxima (Theorem~5). The proof
combines a palindrome symmetry, a transition count that replaces the two
conditions defining a maximum by one, the observation that the relevant
\emph{pairs} of intervals tile once separated by parity, and a small arithmetic
coincidence: the four residues $8x \bmod n$ that arise are always
$\pm1,\pm3,\pm5,\pm7$ in some order, so their squares always sum to $84$.

Nothing in this paper concerns prime numbers.
\end{abstract}

\maketitle

% running head: short title only, no author name on any page
\markboth{An exact histogram for a quadratic staircase}%
        {An exact histogram for a quadratic staircase}

\begin{center}
\emph{Preprint. Not submitted. Circulated for comment.}
\end{center}

\bigskip

\section{The object}

Throughout, $n \ge 3$ is odd and $j$ runs over $0,1,\dots,n-1$. Set
\[
  F(x) = \fl{\frac{2x^2}{n}}, \qquad W_j = F(j+1)-F(j), \qquad
  r_j = 2j^2 \bmod n .
\]
$F$ is a staircase under the parabola $2x^2/n$; $W_j$ is its increment, and
$r_j$ is the residue left over. The two are linked by $2j^2 = nF(j) + r_j$, but
--- and this is the point of Section~\ref{sec:mult} --- the histogram of $W$ can
be determined without ever computing an $r_j$.

Two elementary remarks fix the scale. The total rise is
\[
  \sum_{j=0}^{n-1} W_j = F(n)-F(0) = 2n,
\]
so the mean increment is exactly $2$; and $W_j$ counts the multiples of $n$ in
the interval
\begin{equation}\label{eq:interval}
  I_j = \bigl(2j^2,\ 2(j+1)^2\bigr], \qquad |I_j| = 4j+2 .
\end{equation}
The intervals $I_0, I_1, \dots$ tile $(0,\infty)$. That is the only structural
fact the histogram needs.

\section{The range}

\begin{theorem}\label{thm:range}
For every odd $n$ and every $0 \le j \le n-1$,
\[
  W_j \in \{0,1,2,3,4\}.
\]
\end{theorem}

\begin{proof}
The two arguments of the floor differ by
\[
  \frac{2(j+1)^2}{n}-\frac{2j^2}{n} \;=\; \frac{4j+2}{n},
\]
which is positive and, for $j \le n-1$, strictly less than $4$. A difference of
floors of two reals differing by less than $4$ lies in $\{0,1,2,3,4\}$.
\end{proof}

\begin{verification}
Zero violations over every odd $n \le 20{,}001$.
\end{verification}

The restriction $j \le n-1$ is real: for $j$ in the range $[rn,(r+1)n)$ the same
computation gives $W_j \in \{4r,\dots,4r+4\}$, so the bound moves up in steps of
four. Everything below concerns $0 \le j \le n-1$, which we call the
\emph{first cycle}.

Two further symmetries, both immediate from~\eqref{eq:interval}, will be used:
\begin{equation}\label{eq:sym}
  W_j + W_{n-1-j} = 4, \qquad W_{(n-1)/2} = 2 .
\end{equation}

\section{The exact multiplicities}\label{sec:mult}

\begin{theorem}\label{thm:mult}
Let $n \ge 3$ be odd and $A = \fl{(n+7)/8}$. Over the first cycle,
\[
  \#\{W{=}0\} = \#\{W{=}4\} = A, \qquad \#\{W{=}2\} = 2A-1, \qquad
  \#\{W{=}1\} = \#\{W{=}3\} = \frac{n+1}{2}-2A,
\]
and no other value occurs.
\end{theorem}

\begin{table}[h]
\centering
\begin{tabular}{r l r}
\toprule
$n$ & $(\#0,\ \#1,\ \#2,\ \#3,\ \#4)$ & $A$ \\
\midrule
$11$     & $(2,\ 2,\ 3,\ 2,\ 2)$                       & $2$ \\
$101$    & $(13,\ 25,\ 25,\ 25,\ 13)$                  & $13$ \\
$1{,}009$  & $(127,\ 251,\ 253,\ 251,\ 127)$           & $127$ \\
$19{,}997$ & $(2500,\ 4999,\ 4999,\ 4999,\ 2500)$      & $2500$ \\
\bottomrule
\end{tabular}
\end{table}

\begin{verification}
Zero exceptions for every odd $n \le 200{,}001$.
\end{verification}

\begin{proof}
By~\eqref{eq:interval}, $W_j$ counts the multiples of $n$ in $I_j$, an interval
of length exactly $4j+2$. Hence:
\begin{itemize}
  \item if $4j+2 < n$, the interval is shorter than $n$ and holds
        \emph{at most one} multiple, so $W_j \in \{0,1\}$;
  \item if $4j+2 \ge n$, it is at least as long and holds
        \emph{at least one}, so $W_j \ge 1$.
\end{itemize}

\emph{The count of zeros.} Consequently $W_j = 0$ forces $j < m := \fl{(n+1)/4}$.
On that range the intervals $I_0,\dots,I_{m-1}$ tile $(0,\,2m^2]$ exactly and
each holds at most one multiple, so the number of \emph{occupied} intervals
equals the number of multiples in the range, namely $\fl{2m^2/n}$. Therefore
\begin{equation}\label{eq:N0}
  N_0 \;=\; m - \fl{\frac{2m^2}{n}} .
\end{equation}
Write $n = 8t+r$ with $r \in \{1,3,5,7\}$; then $m = 2t+\varepsilon$ with
$\varepsilon = 0,1,1,2$ respectively. We claim the quotient in~\eqref{eq:N0} is
$q = t+\varepsilon-1$:

\begin{table}[h]
\centering
\begin{tabular}{c c c c c}
\toprule
$r$ & $m$ & $q$ & $2m^2-nq$ & check \\
\midrule
$1$ & $2t$   & $t-1$ & $7t+1$ & $< 8t+1$ \\
$3$ & $2t+1$ & $t$   & $5t+2$ & $< 8t+3$ \\
$5$ & $2t+1$ & $t$   & $3t+2$ & $< 8t+5$ \\
$7$ & $2t+2$ & $t+1$ & $t+1$  & $< 8t+7$ \\
\bottomrule
\end{tabular}
\end{table}

Each remainder is a positive linear function of $t$ strictly below $n$, so the
quotient is as claimed and
\begin{equation}\label{eq:A}
  N_0 = (2t+\varepsilon)-(t+\varepsilon-1) = t+1 = \fl{\frac{n+7}{8}} = A .
\end{equation}

\emph{The remaining four.} Put $H = (n-1)/2$. For $0 \le j < H$ the interval
$I_j$ has length $4j+2 < 2n$, so $W_j \in \{0,1,2\}$. These $H$ intervals tile
$(0,\,2H^2]$, which contains
\[
  \fl{\frac{2H^2}{n}} = H-1
\]
multiples of $n$. Writing $N_i^- = \#\{0 \le j < H : W_j = i\}$, we have
$N_0^- = A$ together with
\[
  N_0^- + N_1^- + N_2^- = H, \qquad N_1^- + 2N_2^- = H-1,
\]
and subtracting gives $N_2^- = A-1$. The central index $j = H$ has $W_H = 2$
by~\eqref{eq:sym}. Finally the symmetry $W_j + W_{n-1-j} = 4$ pairs $W{=}0$ with
$W{=}4$ and $W{=}1$ with $W{=}3$, while $W{=}2$ pairs with itself. Therefore
\[
  N_4 = N_0 = A, \qquad N_2 = 2(A-1)+1 = 2A-1,
\]
and the remaining $n - (N_0+N_2+N_4)$ indices split equally:
\[
  N_1 = N_3 = \frac{n+1}{2}-2A . \qedhere
\]
\end{proof}

\subsection{Why the count has no error term}

A probabilistic model, assuming $2j^2 \bmod n$ equidistributed, estimates $N_0$
by
\[
  \sum_{j < (n-2)/4} \frac{n-4j-2}{n} \;\approx\; \frac{n}{8},
\]
the correct answer, with deviation never exceeding $0.889$ for odd
$n \le 20{,}001$. Sums of this shape normally carry an error of size
$\sqrt{n}\,\log n$ (character sums) or $n^{1/3}$ (lattice points).

The reason none appears here is that the indicators are not independent: the
intervals tile, and the tiling is exact. The residues $2j^2 \bmod n$ never enter
the argument. Theorem~\ref{thm:maxima} is the statement about those residues,
and it is a good deal harder.

\section{Pairs}

Index the symmetric pairs $\{j,\,n-1-j\}$ by $R = 1,3,\dots,n-2$.

\begin{theorem}\label{thm:pairs}
With
\[
  d_n(R) = \fl{\tfrac12 + \frac{(R+2)^2}{2n}} - \fl{\tfrac12 + \frac{R^2}{2n}},
\]
one has $d_n(R) \in \{0,1,2\}$, and the corresponding pair of increments is
$(W_L, W_R) = (2-d,\ 2+d)$.
\end{theorem}

\begin{proof}
The two arguments of the floor differ by
\[
  \frac{(R+2)^2-R^2}{2n} = \frac{2(R+1)}{n},
\]
which for $1 \le R \le n-2$ lies strictly between $0$ and $2$; hence the
difference of floors is $0$, $1$ or $2$. Substituting the pair parametrisation
into the definition of $W$ gives $W_R = 2+d_n(R)$, and~\eqref{eq:sym} then gives
$W_L = 2-d_n(R)$.
\end{proof}

\begin{verification}
Zero failures over every odd $n < 600$ and every odd $R$ with $1 \le R \le n-2$.
\end{verification}

\begin{theorem}\label{thm:paircensus}
Let $N_d = \#\{R : d_n(R) = d\}$. Then $N_2 = N_0+1$.
\end{theorem}

\begin{proof}
There are $(n-1)/2$ odd values $R = 1,3,\dots,n-2$, and the definition of
$d_n(R)$ telescopes along them:
\[
  \sum_R d_n(R) = \fl{\tfrac12 + \frac{n^2}{2n}} - \fl{\tfrac12 + \frac{1}{2n}}
                = \frac{n+1}{2}.
\]
Hence
\[
  N_0+N_1+N_2 = \frac{n-1}{2}, \qquad N_1+2N_2 = \frac{n+1}{2},
\]
and subtracting gives $N_2-N_0 = 1$.
\end{proof}

\begin{verification}
Zero failures over all odd $n < 2000$.
\end{verification}

\begin{table}[h]
\centering
\begin{tabular}{r r r r}
\toprule
$n$ & $N_0$ & $N_1$ & $N_2$ \\
\midrule
$11$  & $1$  & $2$   & $2$ \\
$13$  & $1$  & $3$   & $2$ \\
$31$  & $3$  & $8$   & $4$ \\
$101$ & $12$ & $25$  & $13$ \\
$499$ & $62$ & $124$ & $63$ \\
\bottomrule
\end{tabular}
\end{table}

As in Theorem~\ref{thm:mult}, the whole census is governed by one integer.

\section{The number of local maxima of \texorpdfstring{$2j^2 \bmod n$}{2j\textasciicircum 2 mod n}}

Theorem~\ref{thm:mult} counts the values of $W$. The companion question concerns
their \emph{arrangement}, and its sharpest form is the number of local maxima of
the residue sequence itself.

\begin{theorem}\label{thm:maxima}
Let $n \ge 51$ be odd and $A = \fl{(n+7)/8}$. Then $r_j = 2j^2 \bmod n$,
$j = 0,\dots,n-1$, has exactly $2A$ interior local maxima.
\end{theorem}

\begin{rangeofstatement}
The proof below needs $n \ge 51$, where the interval structure of Step~4 is in
its generic configuration. Computation shows the conclusion is in fact true for
\emph{every} odd $n$ except five: $n = 3,\,5,\,7,\,9,\,49$. Checked exhaustively
for all odd $n \le 200{,}001$.
\end{rangeofstatement}

Throughout put
\[
  \varepsilon_j = [\,r_{j+1} < r_j\,], \qquad
  q_j = \fl{\frac{4j+2}{n}}, \qquad
  \Phi(x) = \fl{\frac{2x^2}{n}} + \fl{\frac{2(x-1)^2}{n}}, \qquad
  H = \frac{n-1}{2},
\]
so that $\varepsilon_j = W_j - q_j$, and $j$ is a local maximum exactly when
$\varepsilon_{j-1} = 0$ and $\varepsilon_j = 1$.

\subsection{Step 1: the palindrome}

Since $2(n-j)^2 \equiv 2j^2$, we have $r_j = r_{n-j}$ for $1 \le j \le n-1$.
Maxima therefore occur in mirror pairs $j \leftrightarrow n-j$, with no fixed
point because $n$ is odd. Also $r_{H+1} = r_{n-H-1} = r_H$, so
\[
  \varepsilon_0 = 0 \quad (\text{since } r_1 = 2 > 0 = r_0), \qquad
  \varepsilon_H = 0,
\]
and no maximum straddles the middle. Hence the total is twice the number of
maxima in $1 \le j \le H$.

\subsection{Step 2: maxima are exactly the mixed pairs}

In any $0$--$1$ word the numbers of ascents and descents differ by the boundary
values; applied to $\varepsilon_0,\dots,\varepsilon_H$ with both ends $0$, the
two are equal. Therefore
\[
  \#\{j \le H : \varepsilon_{j-1}=0,\ \varepsilon_j=1\}
  = \tfrac12\,\#\{j \le H : \varepsilon_{j-1} \ne \varepsilon_j\},
\]
and combining with Step~1,
\[
  \#\{\text{local maxima}\}
  \;=\; \#\{\,1 \le j \le H \;:\; \varepsilon_{j-1} \ne \varepsilon_j\,\}.
\]

\subsection{Step 3: a mixed pair is one condition on one interval}

Since $\varepsilon_{j-1}+\varepsilon_j = (W_{j-1}+W_j) - (q_{j-1}+q_j)$, the
pair at $j$ is mixed if and only if
\[
  P_j := \bigl(2(j-1)^2,\ 2(j+1)^2\bigr]
  \quad\text{contains exactly } q_{j-1}+q_j+1 \text{ multiples of } n.
\]
$P_j$ has length exactly $8j$. No residue occurs in this condition.

\subsection{Step 4: the pair intervals tile, by parity}

Consecutive $P_j$ overlap, but each parity subfamily tiles exactly:
\[
  P_{2s+1} = \bigl(8s^2,\ 8(s+1)^2\bigr], \qquad
  P_{2s}   = \bigl(2(2s-1)^2,\ 2(2s+1)^2\bigr],
\]
of lengths $8(2s+1)$ and $16s$. Since $8j < 4n$ on $j \le H$, each $P_j$ carries
$L_j$ or $L_j+1$ multiples with $L_j = \fl{8j/n} \in \{0,1,2,3\}$, and over any
block of consecutive $j$ of one parity the total telescopes. Set
$\mathrm{off}_j = (q_{j-1}+q_j)-L_j$; then $j$ is mixed if and only if $P_j$
carries the \emph{larger} count when $\mathrm{off}_j = 0$, and the
\emph{smaller} when $\mathrm{off}_j = -1$.

Both $L_j$ and $\mathrm{off}_j$ are explicit step functions of $j$, with jumps
only at
\[
  1 < \ell_1 = \left\lceil \tfrac{n}{8}\right\rceil
    \le \beta_1 = \left\lceil \tfrac{n-2}{4}\right\rceil
    \le \ell_2 = \left\lceil \tfrac{n}{4}\right\rceil
    \le \beta_2 = \left\lceil \tfrac{n+2}{4}\right\rceil
    < \ell_3 = \left\lceil \tfrac{3n}{8}\right\rceil < H < H+1,
\]
giving \emph{seven} intervals with
$(L,\mathrm{off}) = (0,0),(1,-1),(1,0),(2,-1),(2,0),(3,-1),(3,0)$ in order. The
last is the single point $j = H$, where $4H+2 = 2n$ forces $q_H = 2$. Combining
the two parities, an interval $[\alpha,\beta)$ has $\beta-\alpha$ tiles whose
multiples total $\Phi(\beta)-\Phi(\alpha)$, independently of which parity starts
it. Its contribution to the mixed count is $G-LC$ if $\mathrm{off}=0$ and
$(L+1)C-G$ if $\mathrm{off}=-1$, with $G = \Phi(\beta)-\Phi(\alpha)$ and
$C = \beta-\alpha$.

Summing the seven intervals and telescoping the $\Phi$'s (note $\Phi(1)=0$):
\begin{equation}\label{eq:mixed}
  \#\text{mixed} = 2\bigl[\Phi(\ell_1)-\Phi(\beta_1)+\Phi(\ell_2)-\Phi(\beta_2)
                    +\Phi(\ell_3)-\Phi(H)\bigr] + \Phi(H+1) + \Sigma_C,
\end{equation}
\[
  \Sigma_C = 2(\beta_1-\ell_1)-(\ell_2-\beta_1)+3(\beta_2-\ell_2)
             -2(\ell_3-\beta_2)+4(H-\ell_3)-3 .
\]

\begin{verification}
Equation~\eqref{eq:mixed}: zero failures against the direct count for every odd
$9 \le n \le 20{,}001$.
\end{verification}

\subsection{Step 5: evaluating the seven \texorpdfstring{$\Phi$}{Phi}'s}

Two are immediate: $\Phi(H+1) = n-1$ and $\Phi(H) = n-5$. Three more are linear;
writing $n = 8t+r$ with $r \in \{1,3,5,7\}$,
\[
  -\Phi(\beta_1)+\Phi(\ell_2)-\Phi(\beta_2) = -2t + c_r,
  \qquad c_r = 3,\,-1,\,1,\,-3 .
\]
The remaining two are the content, and they are handled together.

\begin{lemma}
For odd $n = 8t+r \ge 51$,
\[
  \Phi\bigl(\lceil n/8\rceil\bigr) + \Phi\bigl(\lceil 3n/8\rceil\bigr)
  = \frac{5n-k_r}{8},
  \qquad k_r = 13,\,7,\,25,\,19 \ \text{ for } r = 1,3,5,7 .
\]
\end{lemma}

\begin{proof}
Put $m = \lceil 3r/8\rceil \in \{1,2,2,3\}$, so $\lceil n/8\rceil = t+1$ and
$\lceil 3n/8\rceil = 3t+m$; the four arguments are
$x \in \{t,\ t+1,\ 3t+m-1,\ 3t+m\}$. Because $8t \equiv -r \pmod n$,
\[
  8x \equiv e_x \pmod n, \qquad e_x = -r,\ 8-r,\ 8m-8-3r,\ 8m-3r,
\]
and in every one of the four cases $\{|e_x|\}$ is a permutation of
$\{1,3,5,7\}$ --- for instance $r=1,m=1$ gives $(-1,7,-3,5)$ and $r=7,m=3$ gives
$(-7,1,-5,3)$. Since $32 \cdot 2x^2 = (8x)^2$,
\[
  32\bigl(2x^2 \bmod n\bigr) \equiv e_x^2 \pmod n,
  \qquad \sum_x e_x^2 = 1+9+25+49 = 84 .
\]
Writing $32(2x^2 \bmod n) = e_x^2 + a_x n$ with $a_x \in [0,32)$ determined by
$e_x^2 + a_x n \equiv 0 \pmod{32}$, a check of the sixteen classes
$n \bmod 32$ gives $\sum_x a_x = 80-4r$ in every case, whence
\[
  \sum_x \bigl(2x^2 \bmod n\bigr) = \frac{84+(80-4r)n}{32}
  = \frac{(20-r)n+21}{8}.
\]
On the other hand $\sum_x 2x^2 = 40t^2+(24m-8)t+(4m^2-4m+4)$ exactly, and
dividing the difference by $n = 8t+r$ gives $5t+s_r$ with
$s_r = -1,1,0,2$, which is $(5n-k_r)/8$.
\end{proof}

\begin{verification}
The Lemma: zero failures for every odd $n \le 300{,}001$ apart from
$n = 3,5,7,9,17,25,49$, where the interval structure of Step~4 degenerates.
\end{verification}

\subsection{Step 6: assembly}

Substituting $\ell_1 = t+1$, $\beta_1 = 2t+\lceil\frac{r-2}{4}\rceil$,
$\ell_2 = 2t+\lceil\frac{r}{4}\rceil$,
$\beta_2 = 2t+\lceil\frac{r+2}{4}\rceil$, $\ell_3 = 3t+m$,
$H = 4t+\frac{r-1}{2}$ into~\eqref{eq:mixed} together with the Lemma gives, in
each of the four residue classes,
\[
  \#\text{mixed} \;=\; 2t+2 \;=\; 2A. \qquad\qquad \blacksquare
\]

\begin{table}[h]
\centering
\begin{tabular}{l c c c c}
\toprule
$n$ & $8t+1$ & $8t+3$ & $8t+5$ & $8t+7$ \\
\midrule
$\#$mixed & $2t+2$ & $2t+2$ & $2t+2$ & $2t+2$ \\
\bottomrule
\end{tabular}
\end{table}

\subsection{Remark: what made it work}

A local maximum is \emph{two} consecutive conditions, while
Theorem~\ref{thm:mult} counts $j$ satisfying \emph{one}; and while the intervals
$I_j$ tile, consecutive $P_j$ overlap. The apparent obstruction dissolves in two
moves. Steps~2 and~3 convert the two conditions into one --- a condition on the
count of multiples in the single interval $P_j$ --- and Step~4 observes that the
$P_j$ tile perfectly once separated by parity. The elementary weight
$\sum_j \varepsilon_j = (n-1)/2$, which follows from $\gcd(4,n)=1$ because the
increments $(4j+2) \bmod n$ then run over every residue, is thereby joined by a
second weight, of pairs rather than of single terms, and the same tiling carries
both.

The one place where genuine arithmetic enters is the Lemma, and it enters
through a small coincidence worth naming: the four numbers $8x \bmod n$ attached
to $\lceil n/8\rceil$ and $\lceil 3n/8\rceil$ are, up to sign, always
$1,3,5,7$ --- so their squares always sum to $84$, whatever $n$ is.

\section{Placement}

Since $r_j = 2j^2 \bmod n$ is $n\{j^2\alpha\}$ with $\alpha = 2/n$
\emph{rational}, the natural comparisons are these.

\medskip\noindent
\textbf{The linear analogue, and what it names.} Replace $j^2$ by $j$. The
difference sequence $u_j = \lfloor (j+1)\alpha\rfloor - \lfloor j\alpha\rfloor$
then takes values in $\{0,1\}$ and is the characteristic \emph{Sturmian} word of
slope $\alpha$; the exact frequencies of its factors are classical, and they are
read off from the three-distance theorem of S\'os, \'Swierczkowski and
Sur\'anyi~\cite{sos} --- see Alessandri and Berth\'e~\cite{ab} for that
connection made explicit. In these terms Theorem~\ref{thm:range} says that $W_j$
is obtained from the same expression with $j$ replaced by $j^2$, over the
alphabet $\{0,1,2,3,4\}$ instead of $\{0,1\}$, and Theorem~\ref{thm:mult} is
its letter census.

\emph{The analogy is one of form only, and it is worth saying how far it fails.}
A Sturmian word is defined by its properties --- a two-letter alphabet, factor
complexity $p(m) = m+1$ for every $m$, aperiodicity, and imbalance at most $1$
--- and not by the displayed formula, which merely characterises it in the
irrational linear case. The word $W$ has none of those properties. It is
periodic; its alphabet has five letters; and its factor complexity, measured on
the first cycle, stabilises for large $n$ at
\[
  p(1),\dots,p(8) \;=\; 5,\ 14,\ 31,\ 56,\ 97,\ 146,\ 219,\ 308
\]
(identical at $n = 4001,\ 10007,\ 20011$), against the Sturmian
$2,3,4,5,6,7,8,9$; its maximum imbalance grows with $n$, measured at
$3,\ 25,\ 39$ for $n = 11,\ 101,\ 1009$, against the Sturmian value $1$. We
found no closed form for the complexity. The point of the comparison is
therefore the name of the object and nothing more --- and it is doubly a
comparison rather than a specialisation, since the linear word at a
\emph{rational} slope is periodic and so is not Sturmian either.

The comparison is worth making because of the one place the two cases part.
For \emph{rational} $\alpha$ the linear word is periodic and its letter
frequencies are immediate. Here the slope is the rational $2/n$, the word is
likewise periodic --- and the multiplicities are not immediate at all: they are
the content of Theorem~\ref{thm:mult}.

Differences of floor functions are of course studied in their own right, and it
is worth saying precisely how the nearest such work differs from the present
object. O'Bryant~\cite{ob} obtains generating functions for sequences
$\lfloor a n/x\rfloor - \lfloor b n/x\rfloor$ and, from them, Beatty-type
partition theorems for $\lfloor 2nx\rfloor - \lfloor nx\rfloor$. That
difference is taken between two dilates of one \emph{linear} sequence at a
common index, and the method requires $x$ irrational, since it runs over the
discontinuities of a step function. Ours is the first difference \emph{along
the index} of a \emph{quadratic} staircase at a \emph{rational} parameter, and
what is counted is the multiplicity of each letter rather than a partition of
the integers. We have not located the quadratic case there or elsewhere, and we
state the difference rather than the resemblance because it is the difference
that leaves the question open. Likewise we do not claim the three-distance
theorem as a template: what the linear case supplies is the name of the family,
not the argument.

\medskip\noindent
\textbf{The quadratic literature runs the other way.} For irrational $\alpha$
the fine-scale statistics of $\{j^2\alpha\}$ --- pair correlation and spacings
--- are known only for almost every $\alpha$ and only as limiting laws;
see Heath-Brown~\cite{hb}. That is a statement about a \emph{continuum} of
$\alpha$ and about a limiting distribution. A theorem of the present kind --- a
finite set of values with explicit multiplicities and no error term --- is
possible only because the orbit here is finite and periodic.

\medskip\noindent
\textbf{The nearest relative} is the literature on spacings of quadratic
residues modulo $q$, for instance Kurlberg~\cite{kurlberg}. The results there
are again distributional; Theorems~\ref{thm:mult}, \ref{thm:paircensus}
and~\ref{thm:maxima} are exact identities for a single explicit sequence.

\medskip
We are not aware of the multiplicities of Theorem~\ref{thm:mult}, the pair
census of Theorem~\ref{thm:paircensus}, or the count of
Theorem~\ref{thm:maxima} in the literature, and would be glad to be corrected.

\appendix

\section{Sample data}

\noindent\textbf{The increments $W_j$, $n = 7$:} $0,\,1,\,1,\,2,\,3,\,3,\,4$.
Here $A = \fl{14/8} = 1$, and Theorem~\ref{thm:mult} predicts
$\#\{0\}=\#\{4\}=1$, $\#\{2\}=2A-1=1$, $\#\{1\}=\#\{3\}=(n+1)/2-2A=2$ ---
exactly what the list shows. (Note that $n=7$ is one of the five small
exceptions to Theorem~\ref{thm:maxima}; the histogram of
Theorem~\ref{thm:mult} has no exceptions at all.)

\medskip\noindent\textbf{The residues $r_j = 2j^2 \bmod 17$:}
$0,\,2,\,8,\,1,\,15,\,16,\,4,\,13,\,9,\,9,\,13,\,4,\,16,\,15,\,1,\,8,\,2$.
Interior local maxima at $j = 2,5,7,10,12,15$ --- six of them, and
$2A = 2\fl{24/8} = 6$.

\medskip\noindent\textbf{The seven intervals of Step~4, $n = 101$}
($t = 12$, $r = 5$, $H = 50$): the boundaries are
\[
  1,\quad \ell_1 = 13,\quad \beta_1 = 25,\quad \ell_2 = 26,\quad
  \beta_2 = 26,\quad \ell_3 = 38,\quad H = 50,\quad H+1 = 51,
\]
so the fourth interval $[\ell_2,\beta_2)$ is empty here and five of the seven
carry indices, with $(L,\mathrm{off}) = (0,0)$ on $[1,13)$, $(1,-1)$ on
$[13,25)$, $(1,0)$ at $j=25$, $(2,0)$ on $[26,38)$, $(3,-1)$ on $[38,50)$ and
$(3,0)$ at $j=50$.

\section{Note on method}

Two details were found by failure rather than by design and are recorded so that
a reader reconstructing the argument does not lose time on them.

\medskip\noindent\textbf{The constant $2$ in the increment.} The author met this
staircase first in a setting where the quantity of interest is $2+W_j$ rather
than $W_j$; dropping the additive constant breaks the formula in $57$ of $78$
tested cases. The constant is an artefact of that setting and not of the
staircase, but it is easy to mislay.

\medskip\noindent\textbf{The seventh interval.} A first version
of~\eqref{eq:mixed} used six intervals and failed for essentially every $n$. The
omission is the single point $j = H$, where $4H+2 = 2n$ forces $q_H = 2$ rather
than $1$; with it restored the identity holds with zero failures on the whole
tested range.

\section*{Reproducibility}

Every count reported above --- the range of $W_j$, the five multiplicities, the
pair census, the seven-interval identity~\eqref{eq:mixed}, the Lemma of Step~5,
and the number of local maxima --- is regenerated by a single script that exits
with a non-zero status if any check fails. It is available from the author on
request.

\section*{Acknowledgement of assistance}

The numerical verification and part of the drafting of this note were carried
out with the assistance of an AI language model; the objects, the questions, the
proofs, and responsibility for every claim are the author's.

\begin{thebibliography}{9}

\bibitem{ab}
P. Alessandri and V. Berth\'e,
\emph{Three distance theorems and combinatorics on words},
Enseign. Math. (2) \textbf{44} (1998), 103--132.

\bibitem{hb}
D.~R. Heath-Brown,
\emph{Pair correlation for fractional parts of $\alpha n^2$},
Math. Proc. Cambridge Philos. Soc. \textbf{148} (2010), 385--407.

\bibitem{kurlberg}
P. Kurlberg,
\emph{The distribution of spacings between quadratic residues, II},
Israel J. Math. \textbf{120} (2000), part A, 205--224.

\bibitem{ob}
K. O'Bryant,
\emph{A generating function technique for Beatty sequences and other step
sequences},
J. Number Theory \textbf{94} (2002), no.~2, 299--319.

\bibitem{sos}
V.~T. S\'os,
\emph{On the distribution mod 1 of the sequence $n\alpha$},
Ann. Univ. Sci. Budapest, E\"otv\"os Sect. Math. \textbf{1} (1958), 127--134.

\end{thebibliography}

\end{document}
